\documentclass{article}
\usepackage{amsmath, amssymb, amsthm}
\usepackage{hyperref}
\usepackage{geometry}
\geometry{margin=1in}

\title{Erd\H{o}s Problem \#835}
\date{Last updated: 2025-08-31}
\author{Source: erdosproblems.com}

\begin{document}
\maketitle

\noindent\textbf{Status:} VERIFIABLE \quad \textbf{No prize}

\noindent\textbf{Tags:} graph theory, hypergraphs

\medskip

\noindent\textbf{Problem Statement:}

Does there exist a $k&#62;2$ such that the $k$-sized subsets of $\{1,\ldots,2k\}$ can be coloured with $k+1$ colours such that for every $A\subset \{1,\ldots,2k\}$ with $\lvert A\rvert=k+1$ all $k+1$ colours appear among the $k$-sized subsets of $A$?




A problem of Erd\H{o}s and Rosenfeld. This is trivially possible for $k=2$. They were not sure about $k=6$. This is equivalent to asking whether there exists $k&gt;2$ such that the chromatic number of the Johnson graph $J(2k,k)$ is $k+1$ (it is always at least $k+1$ and at most $2k$). The chromatic numbers listed at this website show that this is false for $3\leq k\leq 8$. Ma and Tang have proved that the chromatic number of $J(2k,k)$ is $&gt;k+1$ for all $k&gt;2$ not of the form $p-1$ for prime $p$.

\bigskip
\noindent\small{Source: \url{https://www.erdosproblems.com/835}}
\end{document}
